% P03.1 isolated scalar-defect nonremovability proof, live integration.

\begin{proof}[\Cref{thm:p031-isolated-scalar-defect-nonremovability}]
Let $c$ be an isolated point of $\mathcal B_g$.  Choose an open
neighborhood $U_0$ of $c$ such that
\[
                         U_0\cap\mathcal B_g=\{c\}.
\]
The strict inequality $h_g(c)<q_g(c)$ allows us to choose a real number
$\alpha$ with
\[
                         h_g(c)<\alpha<q_g(c).
\]
This also covers $h_g(c)=-\infty$, while the right-hand strict inequality
shows that $q_g(c)$ is finite.  The function $h_g$ is subharmonic by
\Cref{thm:p031-curvewise-scalar-defect-localization}, hence upper
semicontinuous.  After shrinking to an open neighborhood $U\subseteq U_0$
of $c$, we therefore have $h_g<\alpha$ on $U$.

For $d\in U\setminus\{c\}$, the definition of $U_0$, together with
the inequality $h_g\leq q_g$, gives $q_g(d)=h_g(d)$.  It follows that
\[
                    q_g(d)<\alpha<q_g(c)
                    \qquad(d\in U\setminus\{c\}).
\]
Thus $q_g$ has a strict local maximum at $c$.  If it were subharmonic
on some neighborhood of $c$, the one-dimensional maximum principle
\cite[Proposition~3.1.11]{Thu05} would make it locally constant near $c$.
This contradicts the displayed strict inequality, because a smooth analytic
curve has no isolated points: otherwise connectedness would make $C$ a
singleton, contrary to its being a curve.  Hence $c\in\mathcal F_g$.

Suppose now that $\mathcal B_g$ is closed and discrete.  The result just
proved gives $\mathcal B_g\subseteq\mathcal F_g$, while
\Cref{thm:p031-curvewise-scalar-defect-localization} gives
\[
                \mathcal F_g\subseteq\overline{\mathcal B_g}
                =\mathcal B_g.
\]
This proves \eqref{eq:p031-closed-discrete-scalar-defect-exact-failure}.
If $q_g$ is subharmonic, then $\mathcal F_g=\varnothing$, so the
identity forces $\mathcal B_g=\varnothing$.  Conversely, if
$\mathcal B_g=\varnothing$, then the preceding localization theorem gives
subharmonicity of $q_g$ on all of $C$.  A locally finite subset of the
Hausdorff curve $C$ is closed and discrete, which proves the last
assertion.
\end{proof}

\begin{proof}[\Cref{cor:p031-psh-scalar-defect-no-isolated-points}]
If $u$ were classically plurisubharmonic, its pullback
\[
                         q_g=u\circ\pi_{K'/K}\circ g
\]
would be subharmonic on every extended smooth strict test curve by the
definition of classical plurisubharmonicity.  An isolated point of
$\mathcal B_g$ contradicts
\Cref{thm:p031-isolated-scalar-defect-nonremovability}.  This proves both
equivalent formulations.  The final assertion follows from the closed
discrete and locally finite cases of that theorem.
\end{proof}
